% GENERATED FILE. Edit canonical lessons, then run npm run generate:books.
% canonical-source-hash: fnv1a64:3b598441f998ef37
% canonical-lessons: BN-C450-tuchchho, BN-C450-upohas, BN-C450-opoman-kora, BN-C450-khota-deoya, BN-C450-heyo-kora

\chapter{Hurt Feelings}
\label{ch:bn-hurt-feelings}
\hlchaptermodality{450}

\begin{chapteropening}
\textbf{By the end of this chapter:} \emph{I can say that someone insulted, mocked or belittled me.}

Complete the last lesson of chapter 450: I can say that someone insulted, mocked or belittled me.
\end{chapteropening}

\section[\texorpdfstring{\bn{তুচ্ছ} (tuchchhô)}{তুচ্ছ (tuchchhô)}]{\bn{তুচ্ছ} (tuchchhô) — trifling; petty}

\label{lesson:BN-C450-tuchchho}



\begin{quote}
Before the new one: say the Bengali for romantic love, then the Bengali for to be devoted to.
\end{quote}

\subsection*{You'll want to know: \bn{তুচ্ছ}}
\textbf{\bn{তুচ্ছ}} — \emph{tuchchhô} — \textquotedblleft{}trifling; petty\textquotedblright{}.

\begin{quote}
\bn{একটা} \bn{তুচ্ছ} \bn{ব্যাপার} \bn{নিয়ে} \bn{দুই} \bn{ভাই} \bn{ঝগড়া} \bn{করল।}
\end{quote}

\begin{quote}
\emph{ekṭā tuchchhô bæpār niye dui bhāi jhôgṛā kôrlô.}
\end{quote}

\begin{quote}
The two brothers quarrelled over a trifling matter.
\end{quote}

\bn{তুচ্ছ} \bn{করা}, to brush aside, to treat as nothing.

\begin{cousinweb}
\textbf{\bn{তুচ্ছ}} (\emph{tuchchhô}) is a tatsama, Sanskrit \emph{tuccha}, \textquotedblleft{}empty, vain, small, trifling\textquotedblright{} (Monier-Williams).
\end{cousinweb}

\begin{grammarlens}[title={What you've built}]
One more word for describing what you lived and felt.
\end{grammarlens}

\subsection*{Your turn}
\begin{itemize}
\raggedright
  \item \emph{Say it:} \emph{tuchchhô}
  \item \emph{Say it:} \emph{tuchchhô}, once more
  \item \emph{Say it:} \emph{ekṭā tuchchhô bæpār niye dui bhāi jhôgṛā kôrlô.} — The two brothers quarrelled over a trifling matter.
  \item \emph{Recall:} say the Bengali for romantic love, then the Bengali for to be devoted to, then say \emph{tuchchhô} again
\end{itemize}

\begin{tcolorbox}[breakable,colback=teal!4,colframe=teal!35!black,title={Before you move on}]
What does \bn{তুচ্ছ} mean? (\textquotedblleft{}Trifling; petty\textquotedblright{}.) Say it once more.
\end{tcolorbox}
\section[\texorpdfstring{\bn{উপহাস} (upôhās)}{উপহাস (upôhās)}]{\bn{উপহাস} (upôhās) — mockery; derision}

\label{lesson:BN-C450-upohas}



\begin{quote}
Before the new one: say the Bengali for to be devoted to, then the Bengali for trifling.
\end{quote}

\subsection*{You'll want to know: \bn{উপহাস}}
\textbf{\bn{উপহাস}} — \emph{upôhās} — \textquotedblleft{}mockery; derision\textquotedblright{}.

\begin{quote}
\bn{মীরার} \bn{হাসিতে} \bn{উপহাস} \bn{ছিল।}
\end{quote}

\begin{quote}
\emph{mīrār hāsite upôhās chhilô.}
\end{quote}

\begin{quote}
There was mockery in Mira's laugh.
\end{quote}

\bn{উপহাস} \bn{করা}, to mock, to make fun of.

\begin{cousinweb}
\textbf{\bn{উপহাস}} (\emph{upôhās}) is a tatsama, Sanskrit \emph{upahāsa}, \textquotedblleft{}sport, laughter\textquotedblright{} (CDIAL 2282), on the root \emph{has}, to laugh, the \bn{হাসা} (taught) of Bengali; Bengali uses it for laughter aimed at someone.
\end{cousinweb}

\begin{grammarlens}[title={What you've built}]
One more word for describing what you lived and felt.
\end{grammarlens}

\subsection*{Your turn}
\begin{itemize}
\raggedright
  \item \emph{Say it:} \emph{upôhās}
  \item \emph{Say it:} \emph{upôhās}, once more
  \item \emph{Say it:} \emph{mīrār hāsite upôhās chhilô.} — There was mockery in Mira's laugh.
  \item \emph{Recall:} say the Bengali for to be devoted to, then the Bengali for trifling, then say \emph{upôhās} again
\end{itemize}

\begin{tcolorbox}[breakable,colback=teal!4,colframe=teal!35!black,title={Before you move on}]
What does \bn{উপহাস} mean? (\textquotedblleft{}Mockery; derision\textquotedblright{}.) Say it once more.
\end{tcolorbox}
\section[\texorpdfstring{\bn{অপমান} \bn{করা} (ôpômān kôrā)}{অপমান করা (ôpômān kôrā)}]{\bn{অপমান} \bn{করা} (ôpômān kôrā) — to insult; to humiliate}

\label{lesson:BN-C450-opoman-kora}



\begin{quote}
Before the new one: say the Bengali for trifling, then the Bengali for mockery.
\end{quote}

\subsection*{You'll want to know: \bn{অপমান} \bn{করা}}
\textbf{\bn{অপমান} \bn{করা}} — \emph{ôpômān kôrā} — \textquotedblleft{}to insult; to humiliate\textquotedblright{}.

\begin{quote}
\bn{দোকানদার} \bn{খদ্দেরকে} \bn{অপমান} \bn{করল।}
\end{quote}

\begin{quote}
\emph{dokāndār khôdderke ôpômān kôrlô.}
\end{quote}

\begin{quote}
The shopkeeper insulted the customer.
\end{quote}

Past (he/she …): \bn{তিনি} \bn{কর্মচারীকে} \bn{অপমান} \bn{করলেন}, \emph{tini kôrmôchārīke ôpômān kôrlen}.

\begin{cousinweb}
\textbf{\bn{অপমান} \bn{করা}} (\emph{ôpômān kôrā}) is Sanskrit \emph{apamāna}, \textquotedblleft{}disrespect, contempt, disgrace\textquotedblright{} (Monier-Williams), with \emph{kôrā}, \textquotedblleft{}to do\textquotedblright{} (CDIAL 2814).
\end{cousinweb}

\begin{grammarlens}[title={What you've built}]
One more word for describing what you lived and felt.
\end{grammarlens}

\subsection*{Your turn}
\begin{itemize}
\raggedright
  \item \emph{Say it:} \emph{ôpômān kôrā}
  \item \emph{Say it:} \emph{ôpômān kôrā}, once more
  \item \emph{Say it:} \emph{dokāndār khôdderke ôpômān kôrlô.} — The shopkeeper insulted the customer.
  \item \emph{Recall:} say the Bengali for trifling, then the Bengali for mockery, then say \emph{ôpômān kôrā} again
\end{itemize}

\begin{tcolorbox}[breakable,colback=teal!4,colframe=teal!35!black,title={Before you move on}]
What does \bn{অপমান} \bn{করা} mean? (\textquotedblleft{}To insult; to humiliate\textquotedblright{}.) Say it once more.
\end{tcolorbox}
\section[\texorpdfstring{\bn{খোঁটা} \bn{দেওয়া} (khõṭā deoyā)}{খোঁটা দেওয়া (khõṭā deoyā)}]{\bn{খোঁটা} \bn{দেওয়া} (khõṭā deoyā) — to taunt; to throw something in someone's face}

\label{lesson:BN-C450-khota-deoya}



\begin{quote}
Before the new one: say the Bengali for mockery, then the Bengali for to insult.
\end{quote}

\subsection*{You'll want to know: \bn{খোঁটা} \bn{দেওয়া}}
\textbf{\bn{খোঁটা} \bn{দেওয়া}} — \emph{khõṭā deoyā} — \textquotedblleft{}to taunt; to throw something in someone's face\textquotedblright{}.

\begin{quote}
\bn{প্রতিবেশী} \bn{মহিলাটি} \bn{রোজ} \bn{মীরাকে} \bn{খোঁটা} \bn{দিতেন।}
\end{quote}

\begin{quote}
\emph{prôtibeshī môhilāṭi roj mīrāke khõṭā diten.}
\end{quote}

\begin{quote}
The neighbour woman used to taunt Mira every day.
\end{quote}

Past (he/she …): \bn{তিনি} \bn{বোনকে} \bn{খোঁটা} \bn{দিলেন}, \emph{tini bonke khõṭā dilen}.

\begin{cousinweb}
\textbf{\bn{খোঁটা} \bn{দেওয়া}} (\emph{khõṭā deoyā}) is \emph{khõṭā}, \textquotedblleft{}a jibe\textquotedblright{}, whose origin is not settled, with \emph{deoyā}, \textquotedblleft{}to give\textquotedblright{} (CDIAL 6141).
\end{cousinweb}

\begin{grammarlens}[title={What you've built}]
One more word for describing what you lived and felt.
\end{grammarlens}

\subsection*{Your turn}
\begin{itemize}
\raggedright
  \item \emph{Say it:} \emph{khõṭā deoyā}
  \item \emph{Say it:} \emph{khõṭā deoyā}, once more
  \item \emph{Say it:} \emph{prôtibeshī môhilāṭi roj mīrāke khõṭā diten.} — The neighbour woman used to taunt Mira every day.
  \item \emph{Recall:} say the Bengali for mockery, then the Bengali for to insult, then say \emph{khõṭā deoyā} again
\end{itemize}

\begin{tcolorbox}[breakable,colback=teal!4,colframe=teal!35!black,title={Before you move on}]
What does \bn{খোঁটা} \bn{দেওয়া} mean? (\textquotedblleft{}To taunt; to throw something in someone's face\textquotedblright{}.) Say it once more.
\end{tcolorbox}
\section[\texorpdfstring{\bn{হেয়} \bn{করা} (heyô kôrā)}{হেয় করা (heyô kôrā)}]{\bn{হেয়} \bn{করা} (heyô kôrā) — to belittle; to look down on}

\label{lesson:BN-C450-heyo-kora}



\begin{quote}
Before the new one: say the Bengali for to insult, then the Bengali for to taunt.
\end{quote}

\subsection*{You'll want to know: \bn{হেয়} \bn{করা}}
\textbf{\bn{হেয়} \bn{করা}} — \emph{heyô kôrā} — \textquotedblleft{}to belittle; to look down on\textquotedblright{}.

\begin{quote}
\bn{ধনী} \bn{লোকটা} \bn{গরিবদের} \bn{হেয়} \bn{করত।}
\end{quote}

\begin{quote}
\emph{dhônī lokṭā gôribder heyô kôrtô.}
\end{quote}

\begin{quote}
The rich man used to look down on the poor.
\end{quote}

Past (he/she …): \bn{তিনি} \bn{ছাত্রটিকে} \bn{হেয়} \bn{করলেন}, \emph{tini chhātrôṭike heyô kôrlen}.

\begin{cousinweb}
\textbf{\bn{হেয়} \bn{করা}} (\emph{heyô kôrā}) is Sanskrit \emph{heya}, \textquotedblleft{}to be left, to be rejected or avoided\textquotedblright{} (Monier-Williams), with \emph{kôrā}, \textquotedblleft{}to do\textquotedblright{} (CDIAL 2814).
\end{cousinweb}

\begin{grammarlens}[title={What you've built}]
One more word for describing what you lived and felt.
\end{grammarlens}

\subsection*{Your turn}
\begin{itemize}
\raggedright
  \item \emph{Say it:} \emph{heyô kôrā}
  \item \emph{Say it:} \emph{heyô kôrā}, once more
  \item \emph{Say it:} \emph{dhônī lokṭā gôribder heyô kôrtô.} — The rich man used to look down on the poor.
  \item \emph{Recall:} say the Bengali for to insult, then the Bengali for to taunt, then say \emph{heyô kôrā} again
\end{itemize}

\begin{tcolorbox}[breakable,colback=teal!4,colframe=teal!35!black,title={Before you move on}]
What does \bn{হেয়} \bn{করা} mean? (\textquotedblleft{}To belittle; to look down on\textquotedblright{}.) Say it once more.
\end{tcolorbox}
